\section{Declaração de Uso de Ferramentas de Inteligência Artificial}
\label{sec:declaracao-ia}

Em conformidade com os princípios de integridade acadêmica e transparência científica preconizados no regulamento da avaliação, registra-se que os autores empregaram ferramentas baseadas em Modelos de Linguagem de Grande Porte (LLMs) estritamente como suporte auxiliar à redação técnica, formatação tipográfica em \LaTeX\ e estruturação de tabelas conceituais.

Nenhum resultado analítico, julgamento avaliativo ou inferência estatística foi aceito sem auditoria humana integral. As definições metodológicas, o alinhamento com a família de normas ISO/IEC 25000, as proposições de hipóteses e a modelagem GQM foram rigorosamente validados pelos integrantes da equipe. 

A Tabela~\ref{tab:declaracao-ia} discrimina as ferramentas empregadas por integrante, as atividades executadas e os procedimentos de validação humana adotados.

\begin{table}[htbp]
\centering
\caption{Discriminação do Uso de Ferramentas de Inteligência Artificial}
\label{tab:declaracao-ia}
\footnotesize
\setlength{\tabcolsep}{3pt}
\begin{tabular}{|p{2.1cm}|p{2.4cm}|p{4.0cm}|p{4.4cm}|}
\hline
\textbf{Integrante} & \textbf{Ferramenta} & \textbf{Finalidade / Tarefa} & \textbf{Procedimento de Validação Humana} \\ \hline
Daniel & Claude Sonnet 5.5 / ChatGPT & Auxílio na elaboração da estrutura da Seção Metodologia e scripts de automação Git & Revisão manual do fluxo de fases da ISO/IEC 25040, validação das fórmulas de métricas PSM/CID e adequação às regras do repositório. \\ \hline
Heitor & Claude Sonnet 5.5 & Estruturação inicial do texto de introdução, tabela de partes interessadas e configuração MkDocs & Checagem de conformidade com os objetivos reais do Detran-DF, verificação dos critérios de sucesso por parte interessada e remoção de elementos visuais incompatíveis. \\ \hline
Hugo & Claude Code (Claude Opus 5.5) & Apoio na organização do levantamento público dos canais e na revisão do texto da descrição do produto, do modelo de qualidade e dos diagramas & Levantamento e redação feitos pelo autor; dados conferidos nas telas públicas do portal, nas lojas de aplicativos, no CTB e na ISO/IEC 25010:2023. Limitação da Entrega 1: campos que dependem de login serão conferidos na Fase 3. \\ \hline
Pedro Druck & Claude Sonnet 5.5 & Apoio na estruturação da matriz de priorização de características e síntese da família ISO/IEC 25000 & Verificação dos textos oficiais das metas dos ODS (ONU/IBGE) e conferência pontual das subcaracterísticas presentes na versão ISO/IEC 25010:2023. \\ \hline
Pedro Henrique & Claude Sonnet 5.5 & Elaboração do rascunho dos objetivos no template GQM, formulação das hipóteses e referencial teórico GQM/PSM & Auditoria minuciosa para garantir ausência de termos de usabilidade (proibida pela professora); reestruturação das questões para eliminar métricas disfarçadas; checagem das citações bibliográficas de Basili et al. (1994) e McGarry et al. (2002); validação de consistência dos slides em Beamer. \\ \hline
Philipe & Claude Sonnet 5.5 & Estruturação tabular das fichas de métricas e formulação preliminar dos níveis de pontuação & Conferência manual da viabilidade de coleta em caixa-preta sem acesso ao código-fonte; ajuste dos limiares de julgamento conforme referências técnicas de mercado. \\ \hline
\multicolumn{4}{l|}{\scriptsize Fonte: elaborada pela equipe com base nas declarações individuais de uso de ferramentas de inteligência artificial.} \\ \hline
\end{tabular}
\end{table}